% ---
% id: receiver/component-catalogue/def:catalogue-entries
% kind: def
% title: "Entries of the component catalogue"
% summary: >-
%   Ten entry types, from the beamsplitter C1 to the generic block C10, each fix
%   a support, a parameter domain and a passive matrix with every phase and
%   every loss the part has written out. A table collects the supports, the domains and the
%   lossless subfamilies.
% prominence: standard
% status: implemented
% uses: []
% mentions: [receiver/def:catalogue, receiver/def:graph, thm:catalogue-admissibility]
% assumes: []
% refs: []
% provenance: newly-derived
% implements: certiqs_key.receiver.catalogue:Component
% lean:
% notes: "The second forms of eq:retarder, eq:hwpqwp (with the value of HWP(pi/8)) and eq:polarizer are identities; the proof of thm:catalogue-admissibility carries out the products that establish them."
% ---

A receiver declared as an optical layout is built from the entries below,
which form the \tref{component catalogue} of
Definition~\ref{def:catalogue}. Each
entry is written the way the part is specified: every phase and every loss
the part has appears as a parameter, and none is dropped because it could
have been absorbed elsewhere.

\begin{definition}\label{def:catalogue-entries}
The \term[The ten entry types C1 to C10 of the component catalogue, each with a support, a parameter domain and a matrix.]{catalogue entries} are the entry types C1--C10 below.
Throughout, a \term[A spatial index carrying the two polarization modes H and V.]{path} is a spatial index $p$ and carries the two
polarization modes $(p,H)$ and $(p,V)$; a matrix written on two paths acts as
$\cdot\otimes\id_2$ on polarization unless stated otherwise, and a matrix
written on polarization acts on one path.

\emph{C1 --- Beamsplitter.} Support: two paths $(p_1,p_2)$, four modes.
Parameters: transmittance $x\in[0,1]$, reflection phase
$\varphi\in[0,2\pi)$, per-output losses $y_1,y_2\in[0,1]$. Core, on the ordered
path pair with rows $=$ outputs:
\begin{equation}\label{eq:bs}
  \Bsp(x,\varphi)=
  \begin{pmatrix}
    \sqrt{x} & \sqrt{1-x}\,e^{i\varphi}\\[2pt]
    \sqrt{1-x}\,e^{-i\varphi} & -\sqrt{x}
  \end{pmatrix},
  \qquad
  \Gmat_{\mathrm{C1}}=
  \Big[\operatorname{diag}\big(\sqrt{1-y_1},\sqrt{1-y_2}\big)\,
       \Bsp(x,\varphi)\Big]\otimes\id_2 .
\end{equation}
A \term[A beamsplitter that declares its transmittance and reflection phase separately for the H and the V modes.]{polarization-dependent beamsplitter} declares $(x_H,\varphi_H)$ and
$(x_V,\varphi_V)$ separately, with the per-output losses $y_1,y_2$ of C1, and
is
$\operatorname{diag}\big(\sqrt{1-y_1},\sqrt{1-y_2},\sqrt{1-y_1},\sqrt{1-y_2}\big)
\big[\Bsp(x_H,\varphi_H)\oplus\Bsp(x_V,\varphi_V)\big]$ in the reordered basis
$(p_1H,p_2H,p_1V,p_2V)$; this is the matrix of entry C2 at
$(t_H,\varphi_H,t_V,\varphi_V)=(x_H,\varphi_H,x_V,\varphi_V)$ written in that
basis.

\emph{C2 --- Polarizing beamsplitter.} Support: two paths, four modes.
Parameters: $t_H,t_V\in[0,1]$ (ideal $t_H=1$, $t_V=0$), phases
$\varphi_H,\varphi_V\in[0,2\pi)$, losses $y_1,y_2\in[0,1]$. In the ordering
$(p_1H,p_1V,p_2H,p_2V)$,
\begin{equation}\label{eq:pbs}
  \Gmat_{\mathrm{C2}}=
  \operatorname{diag}\big(\sqrt{1-y_1},\sqrt{1-y_1},
                          \sqrt{1-y_2},\sqrt{1-y_2}\big)\;
  \Xi^\dagger\big[\Bsp(t_H,\varphi_H)\oplus\Bsp(t_V,\varphi_V)\big]\,\Xi ,
\end{equation}
where $\Xi$ is the permutation from $(p_1H,p_1V,p_2H,p_2V)$ to
$(p_1H,p_2H,p_1V,p_2V)$.

\emph{C3 --- \term[Entry C3: a wave plate with a retardance, a fast-axis angle and a scalar loss, with half-wave and quarter-wave presets.]{Retarder}, with half- and quarter-wave presets.} Support: the two
polarization modes of one path. Parameters: retardance $\delta\in[0,2\pi)$,
fast-axis angle $\theta\in[0,\pi)$, scalar loss $y\in[0,1]$. With
$\Rot(\theta)=\left(\begin{smallmatrix}\cos\theta&-\sin\theta\\
\sin\theta&\cos\theta\end{smallmatrix}\right)$,
\begin{equation}\label{eq:retarder}
  \Wplate(\delta,\theta)
  \coloneqq e^{-i\delta/2}\,\Rot(\theta)\operatorname{diag}(1,e^{i\delta})
    \Rot(\theta)^{\transpose}
  = e^{-i\delta/2}
  \begin{pmatrix}
    \cos^2\theta+e^{i\delta}\sin^2\theta &
    (1-e^{i\delta})\sin\theta\cos\theta\\[2pt]
    (1-e^{i\delta})\sin\theta\cos\theta &
    \sin^2\theta+e^{i\delta}\cos^2\theta
  \end{pmatrix},
\end{equation}
and $\Gmat_{\mathrm{C3}}=\sqrt{1-y}\,\Wplate(\delta,\theta)$.
The presets are
\begin{equation}\label{eq:hwpqwp}
  \mathrm{HWP}(\theta)\coloneqq\Wplate(\pi,\theta)
  =-i\begin{pmatrix}\cos2\theta&\sin2\theta\\
                    \sin2\theta&-\cos2\theta\end{pmatrix},
  \qquad
  \mathrm{QWP}(\theta)\coloneqq\Wplate(\tfrac{\pi}{2},\theta),
\end{equation}
so that $\mathrm{HWP}(\pi/8)=-\tfrac{i}{\sqrt2}
\left(\begin{smallmatrix}1&1\\1&-1\end{smallmatrix}\right)$.

\emph{C4 --- Circular retarder, or polarization \term[Entry C4: the circular retarder, a real rotation of the two polarization modes of one path.]{rotator}.} Support: the two
polarization modes of one path. Parameter: rotation $\rho\in[0,2\pi)$.
$\Gmat_{\mathrm{C4}}=\Rot(\rho)$.

\emph{C5 --- Mirror or fold.} Support: one path. Parameters: reflectivities
$R_H,R_V\in[0,1]$, phases $\varphi_H,\varphi_V\in[0,2\pi)$, handedness flip
$\mathrm{fl}\in\{\mathrm{false},\mathrm{true}\}$ (default $\mathrm{false}$).
\begin{equation}\label{eq:mirror}
  \Gmat_{\mathrm{C5}}=\operatorname{diag}\big(\sqrt{R_H},\sqrt{R_V}\big)\,
  \operatorname{diag}\big(e^{i\varphi_H},e^{i\varphi_V}\big)\,
  F,\qquad
  F=\begin{cases}\operatorname{diag}(1,-1)&\mathrm{fl}=\mathrm{true}\\
                 \id_2&\text{otherwise.}\end{cases}
\end{equation}

\emph{C6 --- Phase shifter.} Support: any declared mode subset $\Sigma$.
Parameters: $\varphi_m\in[0,2\pi)$ per $m\in\Sigma$.
$\Gmat_{\mathrm{C6}}=\operatorname{diag}(e^{i\varphi_m})_{m\in\Sigma}$.

\emph{C7 --- Loss.} Support: any declared mode subset $\Sigma$. Parameters:
$\eta_m\in[0,1]$ per $m\in\Sigma$, optionally declared as
$L_{\mathrm{dB}}\ge0$ with $\eta=10^{-L_{\mathrm{dB}}/10}$.
$\Gmat_{\mathrm{C7}}=\operatorname{diag}(\sqrt{\eta_m})_{m\in\Sigma}$.

\emph{C8 --- Polarizer.} Support: the two polarization modes of one path.
Parameters: transmission axis $\theta\in[0,\pi)$, pass amplitude $\sqrt{t}$ and
leak amplitude $\sqrt{t_\perp}$ with $t,t_\perp\in[0,1]$. With
$\ket{e_\theta}=\cos\theta\ket{H}+\sin\theta\ket{V}$,
\begin{equation}\label{eq:polarizer}
  \Gmat_{\mathrm{C8}}
  =\sqrt{t}\,\dyad{e_\theta}+\sqrt{t_\perp}\,\dyad{e_\theta^{\perp}}
  =\Rot(\theta)\operatorname{diag}\big(\sqrt{t},\sqrt{t_\perp}\big)
   \Rot(\theta)^{\transpose}.
\end{equation}

\emph{C9 --- Mode routing.} Support: any declared mode subset.
Parameter: a permutation $\pi$ of the modes of $\Sigma$;
$(\Gmat_{\mathrm{C9}})_{ij}=\delta_{i,\pi(j)}$.

\emph{C10 --- \term[Entry C10: an arbitrary matrix on a declared mode subset, admitted after an explicit contractivity check.]{Generic block}.} Support: any declared mode subset $\Sigma$.
Parameter: a matrix $\Gmat_\Sigma\in\mathbb{C}^{\Sigma\times\Sigma}$, admitted
after the explicit check $\Gmat_\Sigma^\dagger\Gmat_\Sigma\preceq\id_\Sigma$.

Table~\ref{tab:catalogue} collects the supports and the parameter domains.

\begin{table}[htbp]
\centering
{\small
\begin{tabular}{@{}l >{\raggedright\arraybackslash}p{0.18\textwidth}
                  >{\raggedright\arraybackslash}p{0.30\textwidth}
                  >{\raggedright\arraybackslash}p{0.32\textwidth}@{}}
\hline
ID & Support & Parameters (domain) &
Lossless exactly when \\
\hline
C1 & 2 paths (4 modes) & $x\in[0,1]$, $\varphi\in[0,2\pi)$,
$y_1,y_2\in[0,1]$ & $y_1=y_2=0$ \\
C2 & 2 paths (4 modes) & $t_H,t_V\in[0,1]$, $\varphi_H,\varphi_V\in[0,2\pi)$,
$y_1,y_2\in[0,1]$ & $y_1=y_2=0$ \\
C3 & 1 path, polarization & $\delta\in[0,2\pi)$, $\theta\in[0,\pi)$,
$y\in[0,1]$ & $y=0$ \\
C4 & 1 path, polarization & $\rho\in[0,2\pi)$ & always \\
C5 & 1 path, polarization & $R_H,R_V\in[0,1]$,
$\varphi_H,\varphi_V\in[0,2\pi)$, $\mathrm{fl}$ & $R_H=R_V=1$ \\
C6 & any $\Sigma$ & $\varphi_m\in[0,2\pi)$ & always \\
C7 & any $\Sigma$ & $\eta_m\in[0,1]$ & every $\eta_m=1$ \\
C8 & 1 path, polarization & $\theta\in[0,\pi)$, $t,t_\perp\in[0,1]$ &
$t=t_\perp=1$ (then $\Gmat_\Sigma=\id$) \\
C9 & any $\Sigma$ & a permutation of $\Sigma$ & always \\
C10 & any $\Sigma$ & $\Gmat_\Sigma$ with
$\Gmat_\Sigma^\dagger\Gmat_\Sigma\preceq\id$ (checked) &
$\Gmat_\Sigma^\dagger\Gmat_\Sigma=\id$ (checked) \\
\hline
\end{tabular}}
\caption{The catalogue entries. Support shapes are stated in modes; the
parameter domains are the declared domains of
Definition~\ref{def:graph}\,(E1); the last column names the subfamily on
which Theorem~\ref{thm:catalogue-admissibility} proves the entry lossless.}\label{tab:catalogue}
\end{table}
\end{definition}
